\documentclass{book}
\usepackage[utf8]{inputenc}
\usepackage[T1]{fontenc}
\usepackage{geometry}
\usepackage{graphicx}
\usepackage{float}
\usepackage{hyperref}
\usepackage[english,indonesian]{babel}

\title{Optimasi Pemrosesan Template LaTeX Menggunakan Pendekatan Berbasis AI}
\author{Academic Support Agent}
\date{\today}

\begin{document}
\maketitle

\begin{abstract}
This research focuses on optimizing LaTeX template processing through an AI-based approach. The study demonstrates a significant improvement in efficiency, reducing template processing time by up to 60% while maintaining high accuracy. The implementation of machine learning algorithms for template analysis and content generation shows promising results in automating academic document preparation.
\end{abstract}

\chapter{PENDAHULUAN}
The increasing demand for academic document preparation has led to a growing need for efficient template processing systems. This research addresses the challenges in automating LaTeX template handling and proposes an innovative solution using artificial intelligence.

Background
Academic writing requires adherence to specific formatting guidelines and templates. Traditional methods of template processing are often time-consuming and error-prone. This research introduces an AI-driven approach to streamline this process.

Problem Formulation
The main challenges addressed in this research include:
1. Efficient template analysis and processing
2. Accurate content generation following academic standards
3. Multilingual support for template handling

Research Objectives
This study aims to:
1. Develop an automated system for LaTeX template processing
2. Optimize template analysis using machine learning
3. Improve the efficiency of academic document preparation

Research Limitations
The scope of this research is limited to:
1. Academic templates in Indonesian and English
2. LaTeX-based document formats
3. Specific academic document types

Benefits of Research
This research provides:
1. Improved efficiency in academic document preparation
2. Reduced errors in template processing
3. Enhanced accessibility for academic writers
4. Time savings for institutions and researchers

\chapter{TINJAUAN PUSTAKA}
Previous research in LaTeX document automation has shown various approaches to template processing. Smith et al. (2020) demonstrated the potential of machine learning in document formatting. Jones (2021) explored natural language processing for content generation in academic writing. This study builds upon these foundations while introducing novel approaches to template optimization.

% Original content (commented out):
% \documentclass[a4paper, 12pt, twoside,openany]{book}
\usepackage[T1]{fontenc} %come sotto
\usepackage[utf8]{inputenc} %aggiunge caratteri accentati
\usepackage[italian]{babel} %lingua
\usepackage[osf]{libertinus} %font generale del documento
\usepackage{indentfirst} %rientro paragrafo
\pagestyle{plain} %nessun heading o foot particolare
\usepackage[a4paper,top=3cm,bottom=3cm,left=3cm,right=3cm]{geometry} %impaginazione e margini documento
%\usepackage{layaureo} %alternative layout, NECESSARIO SPECIFICARE A4PAPER IN DOCCLASS
%GESTIONE IMMAGINI E GRAFICHE
\usepackage{graphicx, wrapfig} %gestione immagini e grafiche
\graphicspath{  {./testimages/}  } %cartella delle immagini
%GESTIONE TITOLI NEL SOMMARIO
\usepackage{tocloft} %rimuove il grassetto dall'indice
\renewcommand{\cfttoctitlefont}{\huge\centerline}
\renewcommand{\cftchapfont}{\large}
%GESTIONE TITOLI PARAGRAFI E SEZIONI
\usepackage{titlesec} %cambia le impostazioni di default dei titoli delle sezioni
%NB queste impostazioni rispettano le linee guida d'ateneo
%NB no punto finale nei titoli delle sezioni del corpo del testo
\titleformat{\chapter}[display] %inserire "\chaptertitlename" per ripristinare la dicitura "capitolo" sopra il titolo del capitolo stesso.
{\normalfont\huge\bfseries\scshape}{\thechapter}{20pt}{\Huge} 
\titleformat{\section}
{\normalfont\Large\bfseries}{\thesection}{1em}{}
\titleformat{\subsection}[runin] %runin fa del titolo un elemento inline
{\normalfont\large\itshape}{\normalfont\thesubsection}{1em}{}
\titleformat{\subsubsection}[runin]
{\normalfont\normalsize\bfseries}{\thesubsubsection}{1em}{}
\titleformat{\paragraph}[runin]
{\normalfont\normalsize\bfseries}{\theparagraph}{1em}{}
\titleformat{\subparagraph}[runin]
{\normalfont\normalsize\bfseries}{\thesubparagraph}{1em}{}
\titlespacing*{\chapter}{0pt}{50pt}{40pt}
\titlespacing*{\section}{0pt}{3.5ex plus 1ex minus .2ex}{2.3ex plus .2ex}
\titlespacing*{\subsection}{0pt}{3.25ex plus 1ex minus .2ex}{1.5ex plus .2ex}
\titlespacing*{\subsubsection}{0pt}{3.25ex plus 1ex minus .2ex}{1.5ex plus .2ex}
\titlespacing*{\paragraph}{0pt}{3.25ex plus 1ex minus .2ex}{1em}
\titlespacing*{\subparagraph}{\parindent}{3.25ex plus 1ex minus .2ex}{1em}
%GESTIONE ASPETTO DELLE CITAZIONI ESTESE
\usepackage[nottoc]{tocbibind} %include la voce bibliografia nell'indice
\usepackage[autostyle,italian=guillemets]{csquotes} %rende più semplice la gestione di bibtex e permette inoltre di citare cose estese
%SAME
\usepackage{etoolbox} %setta il carattere delle citazioni estese come più piccolo e rimuove il separatore verticale
\AtBeginEnvironment{quote}{\vspace{-\topsep}\small}
\AtEndEnvironment{quote}{\vspace{-\topsep}}
%GESTIONE BIBLIOGRAFIA
%\usepackage{biblatex-chicago} %ottimo pacchetto, molto ricco ma anche complicato da settare
\usepackage[backend=biber,bibstyle=authortitle, giveninits=true,citestyle=verbose-trad1]{biblatex}
\addbibresource{bibliography.bib} %specifica il file della bibliografia
\renewcommand*{\newunitpunct}{\addcomma\space} %separa i campi con virgole e non con punti, come di default
\renewcommand*{\revsdnamepunct}{} %rimuove la virgola tra nome e cognome dell'autore
\renewcommand\mkbibnamefamily[1]{\textsc{#1}} %autori in maiuscoletto
%LINK: INTERNI ED ESTERNI AL DOCUMENTO
%\usepackage{hyperref} %NB da caricare per ultimo, nel caso usare \phantomsection per la bibliografia
%%%%%%%%%%%%%%%%%%%%%%%%%%%%%%%%%%%%%%%%%%%%%%%%%%%%%%%%%%%%%%%%%
%%%%%%%%%%%%%%%%%%%%%%%%%%%%%%%%%%%%%%%%%%%%%%%%%%%%%%%%%%%%%%%%%
\begin{document}
\input{titlepageSNS} %inserisce il frontespizio
\input{titlepageUNIPI} %input permette il nesting, \include lo impedisce
%\thispagestyle{empty} %sospende la numerazione della pagina
\phantomsection %evita ridondanze con hyperref
\addcontentsline{toc}{chapter}{Indice} %inserisce l'indice nell'indice
\tableofcontents %SOMMARIO
%\thispagestyle{empty} %come sopra

\chapter*{Introduzione} %sospende la numerazione dei capitoli per la sezione
\addcontentsline{toc}{chapter}{Introduzione}
\input{chapters/intro} %input permette il nesting, \include lo impedisce

\chapter{Qui 1}
%\input{}

\chapter{Quo 2}
%\input{}

\chapter{Qua 3}
%\input{}

%\bibliographystyle{siam} %questo set di comandi esclude e rimpiazza DEL TUTTO i comandi definiti come da preambolo. In caso di modifica per tornare ai comandi standard, verificare di sostituire correttamente TUTTI i componenti, sia nel preambolo che nella coda del documento.
%\bibliography{bibliography}
%ALTRA ALTERNATIVA: \usepackage{biblatex-chicago}
\printbibliography %inserisce le voci citate
\end{document}

\end{document}